\begin{abstract}
Fifteen paired pressure--temperature observations in the data distinguish a visible boiling threshold from an equilibrium coexistence curve. The five sustained-boiling observations have an RMS deviation of 0.25~K from IAPWS saturation, while liquid-labelled observations lie about 5.0~K below it. A Clausius--Clapeyron fit gives an interval-dependent apparent enthalpy of 42.5~kJ/mol. Conditional uncertainty, instrument range and comparison with Martian pressure delimit the inference: the lowest example pressure remains 7.92 times the planetary reference.
\end{abstract}
\begin{IEEEkeywords}
Boiling, saturation pressure, phase equilibrium, nucleation, thermometry.
\end{IEEEkeywords}

\section{Equilibrium and the appearance of bubbles}
The familiar boiling temperature near 100~$^\circ$C belongs to a particular absolute pressure. Equality of liquid and vapour chemical potentials defines a coexistence curve whose slope follows the Clapeyron equation,
\begin{equation}
 \frac{dp_\sigma}{dT}=\frac{\Delta h_{\rm vap}}{T(v_g-v_l)},
 \label{eq:pressure-boiling-clapeyron}
\end{equation}
with specific enthalpy change $\Delta h_{\rm vap}$ and specific volumes $v_g$ and $v_l$. Taking the vapour as ideal, neglecting liquid volume and treating the molar enthalpy as constant gives
\begin{equation}
 \ln(p_\sigma/p_*)=a-\frac{\Delta H_{\rm vap}}{RT}.
 \label{eq:pressure-boiling-log}
\end{equation}
The arbitrary reference pressure $p_*$ makes the logarithm dimensionless. This approximation describes a restricted interval; the numerical comparison below instead uses the temperature-dependent IAPWS saturation relation\cite{iapws1992}.

Surface tension, dissolved gas, nucleation sites and pressure history influence visible bubble formation. Homogeneous and heterogeneous nucleation need not follow a simple equilibrium threshold\cite{blander1975}; isolated bubbles after evacuation are therefore insufficient evidence of coexistence. Observation labels retain this distinction: zero denotes liquid without sustained boiling, one sustained boiling, and two an uncertain judgement. Only label one enters the enthalpy fit.

\section{Paired pressure and liquid temperature}
Figure~\ref{fig:pressure-boiling-apparatus} combines a pressure-rated vessel, a calibrated absolute-pressure instrument and an independent immersed temperature probe. Wet isolation protects dry pressure electronics from water and condensation. The paired readings are transferred manually. Liquid temperature must be measured directly; the barometer chip's temperature cannot substitute for it.

\begin{figure}[t]
\centering
\begin{tikzpicture}[x=1cm,y=1cm]
 \draw[thick,rounded corners] (0.3,0.2) rectangle (3.0,1.8);
 \fill[csblue!12] (0.35,0.25) rectangle (2.95,0.85);
 \draw[csblue] (0.35,0.85)--(2.95,0.85);
 \node at (1.65,0.5) {water};
 \draw[thick] (1.1,2.2)--(1.1,0.6);
 \node[above,align=center] at (1.2,2.2) {liquid\\temperature probe};
 \node[draw,align=center] (pressure) at (6.0,1.4) {absolute-pressure\\instrument};
 \draw[thick] (3,1.4)--(3.6,1.4)--(3.6,0.9)--(4.3,0.9)--(4.3,1.4)--(pressure.west);
 \node[below,align=center] at (3.95,0.8) {wet\\isolation};
 \draw[->,thick] (2.7,1.8)--(2.7,3.0)--(4.9,3.0);
 \node[above] at (4.7,3.0) {controlled evacuation};
\end{tikzpicture}
\caption{Apparatus schematic. Independent instruments measure liquid temperature and chamber absolute pressure while wet isolation protects the pressure connection.}
\label{fig:pressure-boiling-apparatus}
\end{figure}

Expansion and evaporation can cool the liquid while the probe continues to relax. Pairing an immediate pressure with a delayed temperature would combine different states. Settled readings, known probe response and a reversible pressure sequence help distinguish equilibrium from lag and leakage. Gauge pressure requires an atmospheric reference before conversion to absolute pressure; the vessel's operating limits constrain the sequence.

A submerged nucleation site experiences an additional pressure approximately $\rho gh$. Probe depth and bubble location therefore affect comparison with a free-surface saturation curve. Noncondensable gas also separates total pressure from water-vapour partial pressure. Neither contribution should be absorbed into an unexplained thermometer correction.

The fifteen paired observations span 5.04--101.18~kPa and 32.06--96.60~$^\circ$C. Conversion uses $T=T_C+273.15$; temperature is not inferred from pressure. Bounded numerical inversion of the IAPWS saturation relation\cite{iapws1992} supplies $T_\sigma(p)$ for comparison, independently of the observation label.

\begin{figure}[t]
\centering
% Generated by prepare_drafts.py from the supplied CSV files.
\begin{tikzpicture}
\begin{groupplot}[group style={group size=1 by 1,vertical sep=1.15cm},
width=0.98\columnwidth,height=4.0cm,
grid=major,enlarge x limits=false,
legend style={font=\scriptsize,at={(0.5,1.03)},anchor=south,
legend columns=4,draw=none,fill=none,column sep=4pt},
title style={font=\scriptsize,at={(0.02,0.98)},anchor=north west,
fill=white,fill opacity=0.9,text opacity=1,inner sep=2pt},
ylabel style={font=\footnotesize},
xlabel style={font=\footnotesize},tick label style={font=\scriptsize}]
\nextgroupplot[ylabel={Absolute pressure (kPa)},xlabel={Liquid temperature ($^\circ$C)}]
\addplot[csblue,solid,line width=0.8pt] table[col sep=comma,x=c0,y expr={(\thisrow{c1}-(0))*(1)}] {data/drafts/pressure-boiling-saturation.csv};
\addlegendentry{Saturation}
\addplot[csred,dashed,line width=0.8pt,only marks,mark=o] table[col sep=comma,x=c2,y expr={(\thisrow{c1}-(0))*(1)}] {data/drafts/pressure-boiling.csv};
\addlegendentry{All pairs}
\addplot[csteal,dashdotted,line width=0.8pt,only marks,mark=*] table[col sep=comma,x=c2,y expr={(\thisrow{c9}-(0))*(1)}] {data/drafts/pressure-boiling.csv};
\addlegendentry{Boiling}
\end{groupplot}
\end{tikzpicture}

\caption{Absolute pressure versus liquid temperature. The solid curve gives IAPWS saturation; open markers show all paired observations, and filled markers identify sustained boiling.}
\label{fig:pressure-boiling-data}
\end{figure}

\section{Saturation-temperature residuals and apparent enthalpy}
With $r=T-T_\sigma(p)$, the five sustained-boiling observations have mean residual $-0.18$~K and RMS residual 0.25~K (Table~\ref{tab:pressure-boiling-results}). The liquid-labelled and uncertain groups have mean residuals $-5.00$ and $-0.99$~K. Joining all fifteen pairs into one phase boundary would erase this distinction between states and judgement confidence.

\begin{table}[t]
\caption{Sustained-boiling observations, residuals $r=T-T_\sigma(p)$ and propagated standard uncertainties $u_r$, treating the entered pressure and temperature uncertainties as independent standard uncertainties.}
\label{tab:pressure-boiling-results}
\centering\footnotesize
\begin{tabular}{@{}rrrr@{}}
\toprule
$p$ (kPa)&$T_C$ ($^\circ$C)&$r$ (K)&$u_r$ (K)\\
\midrule
90.04&96.60&$-0.10$&0.20\\
60.00&85.62&$-0.30$&0.20\\
39.92&75.35&$-0.46$&0.21\\
25.14&65.09&0.01&0.22\\
10.13&46.00&$-0.05$&0.28\\
\bottomrule
\end{tabular}
\end{table}

Unweighted least squares of $\ln(p/1\unit{kPa})$ against $1/T$ for these five observations gives $\Delta H_{\rm app}=42.5\unit{kJ\,mol^{-1}}$ through Eq.~\eqref{eq:pressure-boiling-log}. This interval-dependent value summarizes the thermal dependence of vapour pressure; it does not resolve latent-heat variation or vapour nonideality. Regression scatter alone cannot establish its thermodynamic uncertainty.

At 39.92~kPa, the liquid temperature is 75.35~$^\circ$C rather than the reference 75.81~$^\circ$C. The discrepancy is small on the graph but material relative to the entered uncertainties. Its interpretation depends on probe response, pressure referencing and the boiling judgement.

\section{Uncertainty and the accessible inference}
Interpreting the entered 0.1~kPa and 0.2~K as independent standard uncertainties gives
\begin{equation}
 u_r^2=u_T^2+\left(\frac{u_p}{dp_\sigma/dT}\right)^2.
 \label{eq:pressure-boiling-uncertainty}
\end{equation}
This conditional propagation excludes lag, common calibration bias and judgement uncertainty. The largest normalized residual is 2.20. At lower pressure, the smaller saturation-curve slope gives pressure uncertainty a larger temperature-equivalent contribution, explaining the rise of $u_r$ towards 10~kPa. A common thermometer offset would not diminish through averaging five observations.

\section{Instrument range and the planetary comparison}
The lowest example pressure, 5.039~kPa, remains 7.92 times NASA's 0.636-kPa Martian reference at mean radius; the reported seasonal interval is about 0.40--0.87~kPa\cite{nasamarsfacts}. Water's triple point is approximately 0.611657~kPa at 273.16~K. IAPWS inversion at 0.636~kPa gives a conditional saturation temperature of about 0.55~$^\circ$C. The stable-liquid relation must not be extrapolated below the triple-point pressure.

The LPS22HB and BMP280 lower limits, 26 and 30~kPa\cite{lps22hb,bmp280}, are about 41 and 47 times the Martian reference. Five example states, including two sustained-boiling observations, lie below both limits. Covering even the full example interval therefore requires an appropriate lower-range instrument and wet-compatible connection.

Thermodynamically favourable coordinates alone do not establish persistent Martian liquid water: ice and heat availability, atmospheric undersaturation and evaporation matter, while salts change the phase relations\cite{haberle2001}. The observations establish the pressure dependence of the visible boiling threshold over the measured interval.
